% !TEX root = ../../ctfp-print.tex

\lettrine[lhang=0.17]{我}{们已经见过几种幺半群}（monoid）的表述形式：作为集合、单对象范畴，以及在幺半范畴中的一个对象。这个简单的概念还能衍生出多少新内容？

让我们尝试探索。采用如下定义：幺半群是一个带有两个函数的集合 $m$：
\begin{align*}
  \mu  & \Colon m\times{}m \to m \\
  \eta & \Colon 1 \to m
\end{align*}
其中 $1$ 是 $\Set$ 中的终对象（终端对象）——即单元素集合。
第一个函数定义乘法（接收两个元素返回它们的乘积），第二个函数从 $m$ 中选取单位元。并非任意选择这两个函数都能构成幺半群，还需要满足结合律和单位元律。但暂时忽略这些条件，仅考虑「潜在的幺半群」。这对函数属于两个函数集的笛卡尔积。我们知道这些函数集可以表示为指数对象：
\begin{align*}
  \mu  & \in m^{m\times{}m} \\
  \eta & \in m^1
\end{align*}
这两个集合的笛卡尔积为：
$$m^{m\times{}m}\times{}m^1$$
利用一些中学代数知识（这在每个笛卡尔闭范畴中都成立），可以将其重写为：
$$m^{m\times{}m + 1}$$
这里的 $+$ 表示 $\Set$ 中的余积（coproduct）。我们刚刚将一对函数替换为单个函数——即以下集合的元素：
$$m\times{}m + 1 \to m$$
该函数集的任何元素都是潜在的幺半群。

这种表述的美妙之处在于它能引出有趣的推广。例如，如何用这种语言描述群？群是带有额外逆元函数的幺半群。后者是形如 $m \to m$ 的函数。以整数为例，加法作为二元运算，零作为单位元，取负作为逆元运算，即可构成群。定义群时可以从三元组函数开始：
\begin{align*}
  m\times{}m \to m \\
  m \to m          \\
  1 \to m
\end{align*}
同样地，我们可以将所有这样的三元组合并成一个函数集：
$$m\times{}m + m + 1 \to m$$
我们从一个二元运算（加法）、一个一元运算（取反）和一个零元运算（单位元——此处为零）出发，将它们合并为一个函数。所有具有此签名的函数都定义了潜在的群。

我们可以继续扩展。例如定义环时需要增加另一个二元运算和一个零元运算，依此类推。每次得到的函数类型左侧都是若干幂次（可能包含零次幂——终对象）的和，右侧则是集合自身。

现在我们可以对推广进行更疯狂的探索。首先，可以将集合替换为对象，函数替换为态射。我们可以将 $n$ 元运算定义为从 $n$ 元积到对象的态射。这意味着我们需要支持有限积的范畴。对于零元运算，我们需要存在终对象。因此需要笛卡尔范畴。为了组合这些运算，我们需要指数对象，即要求是笛卡尔闭范畴。最后，为了完成我们的代数操作，还需要余积。

另一种思路是完全忘记公式的推导过程，专注于最终结果。我们态射左侧的积之和定义了一个自函子（endofunctor）。如果我们选择任意自函子 $F$ 会怎样？在这种情况下，我们无需对范畴施加任何约束。这样得到的对象称为 $F$-代数（F-algebra）。

$F$-代数由三部分组成：一个自函子 $F$、一个对象 $a$ 和一个态射
$$F a \to a$$
该对象常被称为载体（carrier），或在编程语境下的载体\emph{类型}。该态射常被称为求值函数（evaluation function）或结构映射。可以将函子 $F$ 视为生成表达式，而态射则用于求值这些表达式。

以下是Haskell中$F$-代数的定义：

\src{snippet01}
它通过求值函数直接标识代数本身。

在幺半群的例子中，对应的函子为：

\src{snippet02}
这是Haskell中 $1 + a\times{}a$ 的表示（回想
\hyperref[simple-algebraic-data-types]{代数数据结构}）。

环可以用以下函子定义：

\src{snippet03}
对应Haskell代码 $1 + 1 + a\times{}a + a\times{}a + a$。

环的一个例子是整数集合。我们可以选择
\code{Integer} 作为载体类型，并定义求值函数为：

\src{snippet04}
基于相同函子 \code{RingF} 的 $F$-代数还有许多其他实例。例如多项式构成环，方阵也构成环。

可以看出，函子的作用是生成可被代数求值器处理的表达式。目前我们只看到非常简单的表达式。通常我们会关注通过递归定义的更复杂表达式。

\section{递归}

生成任意表达式树的一种方法是在函子定义中将变量 \code{a} 替换为递归结构。例如，环中的任意表达式可以通过如下树状数据结构生成：

\src{snippet05}
我们可以用递归版本替换原始的环求值器：

\src{snippet06}
虽然仍需将所有整数表示为1的累加（不够实用），但在某些场景下足够使用。

但如何用$F$-代数的语言描述表达式树？我们需要某种方式形式化在函子定义中递归地用替换结果填充自由类型变量的过程。想象分步骤进行：首先定义深度为1的树为：

\src{snippet07}
我们在 \code{RingF} 定义中用深度为0的树（由 \code{RingF a} 生成）填充空位。深度为2的树类似地定义为：

\src{snippet08}
也可以写作：

\src{snippet09}
继续这一过程，可以写出符号方程：

\begin{snipv}
type RingF\textsubscript{n+1} a = RingF (RingF\textsubscript{n} a)
\end{snipv}
理论上，无限重复此过程后得到我们的 \code{Expr}。注意 \code{Expr} 不依赖于 \code{a}。初始起点无关紧要，最终总会到达相同结果。这在任意范畴中的任意自函子未必成立，但在 $\Set$ 范畴中性质良好。

当然，这只是直观解释，稍后我会给出更严谨的论述。

将自函子无限次应用会产生一个\newterm{不动点}（fixed point），其定义为：
$$\mathit{Fix}\ f = f\ (\mathit{Fix}\ f)$$
该定义的直觉是：由于我们已对 $f$ 应用了无限次得到 $\mathit{Fix}\ f$，再应用一次不会改变结果。在Haskell中，不动点定义为：

\src{snippet10}
严格来说，若构造器名称与类型名称不同会更具可读性，例如：

\src{snippet11}
但我们将遵循公认记法。构造器 \code{Fix}（或 \code{In}）可视作函数：

\src{snippet12}
还存在一个函数用于剥离一层函子应用：

\src{snippet13}
这两个函数互为逆元。后续我们将使用这些函数。

\section{F-代数范畴}

这是最古老的技巧：每当你想出一种构造新对象的方法时，都要验证它们是否构成一个范畴。不出所料，给定自函子 $F$ 的代数确实构成一个范畴。该范畴的对象就是代数——由载体对象 $a$ 和态射 $F a \to a$ 组成的有序对，两者都来自原始范畴 $\cat{C}$。

为了完善这个图像，我们需要定义F-代数范畴中的态射。一个态射必须将代数 $(a, f)$ 映射到另一个代数 $(b, g)$。我们将它定义为载体间的态射 $m$——即在原始范畴中从 $a$ 到 $b$ 的态射。但并非任意态射都满足条件：我们要求它与两个求值映射兼容。（我们称这种结构保持态射为\newterm{同态}。）以下是F-代数同态的定义方式：首先注意到我们可以将 $m$ 提升为映射
$$F m \Colon F a \to F b$$
然后将其与 $g$ 复合得到 $b$。等价地，我们可以先用 $f$ 将 $F a$ 映射到 $a$，再通过 $m$ 复合。我们要求这两条路径相等：
$$g \circ F m = m \circ f$$

\begin{figure}[H]
  \centering
  \includegraphics[width=0.3\textwidth]{images/alg.png}
\end{figure}

\noindent
容易验证这确实构成一个范畴（提示：原始范畴 $\cat{C}$ 中的恒等态射完全适用，且同态的复合仍是同态）。

若F-代数范畴存在初始对象，则称为\newterm{初始代数}。记该初始代数的载体为 $i$，求值映射为 $j \Colon F i \to i$。实际上，初始代数的求值映射 $j$ 是一个同构。这个结论被称为Lambek定理。证明依赖于初始对象的定义：要求存在从它到任意其他F-代数的唯一同态 $m$。由于 $m$ 是同态，下述图必须交换：

\begin{figure}[H]
  \centering
  \includegraphics[width=0.3\textwidth]{images/alg2.png}
\end{figure}

\noindent
现在构造一个载体为 $F i$ 的代数。此类代数的求值映射必然是从 $F (F i)$ 到 $F i$ 的态射。我们可以通过提升 $j$ 轻松构造这样的求值映射：
$$F j \Colon F (F i) \to F i$$
因为 $(i, j)$ 是初始代数，必然存在从它到 $(F i, F j)$ 的唯一同态 $m$。下述图必须交换：

\begin{figure}[H]
  \centering
  \includegraphics[width=0.3\textwidth]{images/alg3a.png}
\end{figure}

\noindent
但我们还有这个显然交换的图（两条路径完全相同！）：

\begin{figure}[H]
  \centering
  \includegraphics[width=0.3\textwidth]{images/alg3.png}
\end{figure}

\noindent
该图可解释为表明 $j$ 是将 $(F i, F j)$ 映射到 $(i, j)$ 的代数同态。将这两个图拼接得到：

\begin{figure}[H]
  \centering
  \includegraphics[width=0.6\textwidth]{images/alg4.png}
\end{figure}

\noindent
该图表明 $j \circ m$ 是代数同态。特别地，此时两个代数是相同的。更进一步，由于 $(i, j)$ 是初始对象，从它到自身的同态只能有一个，即恒等态射 $\id_i$——我们已知它是代数同态。因此 $j \circ m = \id_i$。利用此事实和左图的交换性可证 $m \circ j = \id_{Fi}$。这表明 $m$ 是 $j$ 的逆，从而 $j$ 是 $F i$ 与 $i$ 之间的同构：
$$F i \cong i$$
也就是说 $i$ 是 $F$ 的不动点。这就是最初直观论证的严格证明。

回到Haskell：我们识别出 $i$ 对应类型 \code{Fix f}，$j$ 对应构造器 \code{Fix}，其逆对应 \code{unFix}。Lambek定理的同构告诉我们，为获得初始代数，我们取函子 $f$ 并将其参数 $a$ 替换为 \code{Fix f}。这也解释了为何不动点不依赖于类型参数 $a$。

\section{自然数}

自然数也可以被定义为一个F-代数(F-algebra)。其起点是一对态射(morphisms)：
\begin{align*}
  zero & \Colon 1 \to N \\
  succ & \Colon N \to N
\end{align*}
第一个态射选取零元素，第二个态射将每个数映射到它的后继。与之前类似，我们可以将这两个映射合并为一个：
$$1 + N \to N$$
左边定义了一个函子(functor)，在Haskell中可以写成：

\src{snippet14}
这个函子的不动点(fixed point)（即它生成的初始代数(initial algebra)）在Haskell中可以编码为：

\src{snippet15}
一个自然数要么是零，要么是另一个数的后继。这被称为自然数的皮亚诺(Peano)表示法。

\section{折叠}

让我们用Haskell记法重写初始性条件。我们将初始代数称为\code{Fix f}，它的求值器是构造函数\code{Fix}。对于该函子上的任意其他代数，存在唯一的态射\code{m}从初始代数映射到它。我们选取一个载体类型为\code{a}、求值器为\code{alg}的代数：

\begin{figure}[H]
  \centering
  \includegraphics[width=0.4\textwidth]{images/alg5.png}
\end{figure}

\noindent
顺便注意，\code{m}本质上是什么：它是不动点的求值器，是整个递归表达式树的求值器。现在我们寻找实现它的通用方式。

根据拉姆贝克(Lambek)定理，构造函数\code{Fix}是一个同构。我们将其逆映射记作\code{unFix}。因此可以反转图示中的箭头得到：

\begin{figure}[H]
  \centering
  \includegraphics[width=0.4\textwidth]{images/alg6.png}
\end{figure}

\noindent
写出该图示的交换条件：

\begin{snip}{haskell}
m = alg . fmap m . unFix
\end{snip}
我们可以将此方程解读为\code{m}的递归定义。对于通过函子\code{f}创建的任何有限树，递归必定会终止。这可以通过观察到\code{fmap m}作用于函子\code{f}的底层结构来验证——换句话说，它作用于原始树的子节点上，而子节点的深度总比原树少一层。

当我们将\code{m}应用于由\code{Fix f}构造的树时会发生以下过程：\code{unFix}的操作剥离构造函数，暴露树的顶层结构。接着我们对顶层节点的所有子节点应用\code{m}，这将产生类型为\code{a}的结果。最后通过非递归求值器\code{alg}组合这些结果。关键在于我们的求值器\code{alg}是一个简单的非递归函数。

由于对任意代数\code{alg}都可以执行此操作，我们应当定义一个高阶函数：以代数为参数，返回我们所需的函数\code{m}。这个高阶函数被称为折叠(cata-)：

\src{snippet16}
来看具体实例。取定义自然数的函子：

\src{snippet17}
选择载体类型为\code{(Int, Int)}，定义代数如下：

\src{snippet18}
容易验证该代数对应的折叠\code{cata fib}将计算斐波那契数列。

一般而言，针对\code{NatF}的代数定义了一个递推关系：当前元素的值由前序元素决定。折叠则计算该序列的第n项。

\section{折叠}

\code{e}的列表是以下函子的初始代数：

\src{snippet19}
事实上，将变量\code{a}替换为递归结果（我们称之为\code{List e}），可得：

\src{snippet20}
列表函子的一个代数会选择特定的载体类型，并定义一个对两个构造器进行模式匹配的函数。该代数对\code{NilF}的取值表示如何求值空列表，对\code{ConsF}的取值则表示如何将当前元素与之前累积的值进行组合。

例如，以下这个代数可用于计算列表长度（载体类型为\code{Int}）：

\src{snippet21}
事实上，由此产生的叠积（catamorphism）\code{cata lenAlg}可以计算列表长度。请注意，这个求值器由两部分组成：(1) 一个接受列表元素和累加器并返回新累加器的函数；(2) 一个初始值（此处为零）。累加器的类型和值的类型由载体类型决定。

将其与传统的Haskell定义对比：

\src{snippet22}
\code{foldr}的两个参数恰好对应代数的两个分量。

让我们尝试另一个例子：

\src{snippet23}
同样地，将其与以下代码对比：

\src{snippet24}
正如所见，\code{foldr}本质上就是列表范畴上叠积（catamorphism）的一个便捷特例。

\section{余代数}

与往常一样，我们有一个对偶构造F-余代数，其中态射的方向被反转：
$$a \to F a$$
给定函子的余代数也构成一个范畴，其同态保持余代数结构。该范畴的终对象$(t, u)$称为终端(或终结)余代数。对于任意其他余代数$(a, f)$，存在唯一的同态$m$使得下图可交换：

\begin{figure}[H]
  \centering
  \includegraphics[width=0.4\textwidth]{images/alg7.png}
\end{figure}

\noindent
终余代数是函子的一个不动点，即态射$u \Colon t \to F t$是一个同构(Lambek定理对于余代数)：
$$F t \cong t$$
终余代数在编程中通常被解释为生成(可能无限的)数据结构或转移系统的配方。

正如catamorphism可用于求值初始代数，anamorphism可用于余代数的余评估：

\src{snippet25}
余代数的经典例子基于某个函子，其不动点是类型为\code{e}元素的无限流。这个函子是：

\src{snippet26}
其固定点为：

\src{snippet27}
对于\code{StreamF e}的余代数是一个函数：它接受类型为\code{a}的种子，产生一个由元素和下一个种子组成的对(\code{StreamF}是对"对"的 fancy 称呼)。

你可以轻松生成产生无限序列的余代数示例，比如平方数列表或倒数序列。

更有趣的例子是生成素数列表的余代数。技巧在于使用无限列表作为载体。我们的初始种子是列表\code{{[}2..{]}}。下一个种子是移除所有2的倍数后的尾部列表，即从3开始的奇数列表。下一步取该列表的尾部并移除所有3的倍数，依此类推。你可能认出这是埃拉托斯特尼筛法的雏形。该余代数通过以下函数实现：

\src{snippet28}
该余代数的anamorphism生成素数列表：

\src{snippet29}
流是无限列表，因此应该可以转换为Haskell列表。为此，我们可以使用相同的函子\code{StreamF}构造代数，并在其上运行catamorphism。例如，这是一个将流转换为列表的catamorphism：

\src{snippet30}
这里，相同不动点同时作为该自函子的初始代数和终余代数。但在任意范畴中并非总是如此。一般来说，自函子可能有多个(或没有)不动点。初始代数称为最小不动点，终余代数称为最大不动点。然而在Haskell中，两者都由相同公式定义且重合。

列表的anamorphism被称为unfold。为创建有限列表，函子被修改为生成\code{Maybe}对：

\src{snippet31}
当值为\code{Nothing}时终止列表生成。

有趣的是，镜头(lens)与余代数有关。镜头可表示为获取函数(getter)和设置函数(setter)的组合：

\src{snippet32}
此处\code{a}通常是包含类型为\code{s}字段的乘积数据类型。获取函数检索该字段值，设置函数用新值替换该字段。这两个函数可合并为一个：

\src{snippet33}
我们可以进一步改写为：

\src{snippet34}
其中我们定义了函子：

\src{snippet35}
注意这不是由和与积构造的简单代数函子，它包含指数$a^s$。

镜头是该函子以\code{a}为载体类型的余代数。我们之前已看到\code{Store s}也是一个余单子(comonad)。事实证明，良好行为的镜头对应兼容余单子结构的余代数。我们将在下一节讨论这一点。

\section{挑战}

\begin{enumerate}
  \tightlist
  \item
        实现一个单变量多项式环的求值函数。可以将多项式表示为按升幂排列的系数列表。例如，$4x^2-1$ 的表示形式为（从零次幂开始）
        \code{{[}-1, 0, 4{]}}。
  \item
        将上述构造推广至多独立变量的多项式，如 $x^2y-3y^3z$。
  \item
        实现 $2\times{}2$ 矩阵环的代数结构。
  \item
        定义一个余代数(coalgebra)，其叠核(anamorphism)生成自然数平方的列表。
  \item
        使用 \code{unfoldr} 函数生成前 $n$ 个素数(primes)的列表。
\end{enumerate}